\section{The all-field multidimensional substitution statement}

For a finite-dimensional space $W=V_1\otimes\cdots\otimes V_b$
and a $j$-plane $L\subset W$, write
\[
 \mathsf P_L=\sum_{i=1}^j a_i\otimes u_i
       \in k^j\otimes V_1\otimes\cdots\otimes V_b,
\]
where $(u_i)$ is any basis of $L$. Its CP border rank is independent
of the basis. All the $b$ residual factors remain separate.

\begin{finding}{R-01}{proper incidence and substitution of $j$ dimensions}
Let $k$ be algebraically closed, let
$T\in A_0\otimes V_1\otimes\cdots\otimes V_b$ be nonzero,
and let $U$ be its $A_0$-flattening image, of dimension $q$.
For $1\leq j\leq q$,
\begin{equation}\label{eq:substitution}
 \BR_k(T)\ \geq\ q-j+
       \min_{L\in\Gr(j,U)}\BR_k(\mathsf P_L).
\end{equation}
The same inequality lower-bounds an original tensor from which $T$
was obtained by grouping factors into $A_0$.
\end{finding}
\begin{proof}
For $s\geq1$ let $Z_{j,s}\subset\Gr(j,W)$ be the locus
$\BR(\mathsf P_L)\leq s$. It is closed. Indeed, on each
standard Grassmannian chart choose the regular echelon frame of
the tautological bundle. The associated tensor is a regular
map into $k^j\otimes W$. The inverse image of its closed
Segre-secant cone is closed in that chart. Changing the frame
acts through $\GL_j$, which preserves the cone, so these local
conditions agree and define the claimed closed locus.

In $\Gr(j,W)\times\Gr(q,W)$ the flag incidence $L\subset U'$
is closed: it is the vanishing of the map from the rank-$j$
tautological subbundle to the rank-$(\dim W-q)$ universal
quotient. Therefore
\[
 E_{q,j,s}=\{U':\text{some }L\subset U',\ \dim L=j,
                         \BR(\mathsf P_L)\leq s\}
\]
is closed. It is the image of this incidence intersected with
$Z_{j,s}\times\Gr(q,W)$ under the projective projection.
The precise properness inputs are EGA~II, Definition~5.4.1,
Proposition~5.4.2 and Theorem~5.5.3(i); see the dictionary.

Replace $A_0$ by the $q$-dimensional support of $T$ in that
factor. An inclusion and a left-inverse projection show that
this replacement preserves its border rank. On the open set
$O$ of tensors whose first flattening has full rank $q$, the
image map $f:O\to\Gr(q,W)$ is regular: its coordinates on
a nonzero-minor chart are the usual echelon-frame coordinates.
Put $r=\BR(T)$; an ordinary flattening gives $r\geq q$.

For an actual rank-at-most-$r$ tensor $T'\in O$, let $S$ be
the span of the pure residual atoms of a length-at-most-$r$
decomposition. Then $U'=f(T')\subset S$ and
$q\leq d:=\dim S\leq r$. Choose a basis of $S$ from those
pure atoms and let $M$ be the span of any $d-q+j$ of them.
The dimension inequality gives $\dim(U'\cap M)\geq j$.
Choose a $j$-plane $L\subset U'\cap M$. Expanding its
basis in the selected atoms expresses $\mathsf P_L$ as a
sum of at most $d-q+j\leq r-q+j$ pure tensors, including
its new first factor. Thus $f(T')\in E_{q,j,r-q+j}$.

The set $f^{-1}(E_{q,j,r-q+j})$ is closed in $O$. Since
$T$ is in the closure in $O$ of the actual rank-at-most-$r$
tensors in $O$, it also belongs to that inverse image. Hence
some $L\subset U$ has $\BR(\mathsf P_L)\leq r-q+j$.
Rearranging proves \eqref{eq:substitution}. Grouping original
factors maps every original pure tensor to a pure tensor in
the grouped format and hence cannot increase border rank;
the last assertion follows.
\end{proof}

The multidimensional three-factor principle is already
Landsberg--Micha{\l}ek, Proposition~2.3\cite{LM2018}.
The proof above records the precise arbitrary residual-order,
arbitrary-algebraically-closed-field form needed here. Its
$j=1$ case is the proper-incidence version of predecessor
GCT\_GINEQ F-13/F-18; no new attribution of the substitution
principle is intended. The new numerical work is the residual
minimum and its certificates.

\paragraph{Geometric form.}
Each $L$ is a point of the Grassmannian with its tautological
rank-$j$ locally free sheaf. Its inclusion into the tensor
bundle is the tautological residual tensor. Border rank forces
the image point $[U]$ into the proper image $E_{q,j,s}$ of a
closed flag incidence. The proof is carried out in each
algebraically closed fiber. It does not assume that an
integral secant ideal commutes with every reduction modulo a
prime, and it uses no Euclidean compactness or Hermitian norm.
